% ============================================================
% 第八章 解析延拓的定义与唯一性
% 内容：定义、恒等定理、函数关系范式、本书三种途径、自然边界、与 ζ 的衔接
% 目的：衔接第7章 Γ 构造与第9章 ζ 的具体延拓
% 编译方式：XeLaTeX（作为 main.tex 的 \input 子文件）
% ============================================================

\chapter{解析延拓的定义与唯一性}\label{ch:analytic_continuation_general}

\section{引言}\label{sec:intro_why_continuation}

许多复变函数最初只由在局部区域（圆盘、半平面等）内收敛的级数或积分定义；在该区域外，同一表达式可能发散。
收敛边界常常只是\textbf{表示式}的限制，而非函数本身不能延拓。
典型例子：几何级数 $\sum z^n$ 仅在 $|z|<1$ 收敛，而 $1/(1-z)$ 在 $\mathbb{C}\setminus\{1\}$ 上全纯；
$\zeta(s)$ 的 Dirichlet 级数在 $\operatorname{Re}(s)\le 1$ 发散，但 $\zeta$ 仍可亚纯延拓到全平面。

\textbf{解析延拓}问的是：给定区域 $D_0$ 上全纯的 $f$，是否存在更大区域 $D\supset D_0$ 及 $D$ 上全纯的 $F$，使 $F|_{D_0}=f$。
若存在，则在 $D$ 连通时延拓唯一（下节）。
本章先给定义与唯一性，再给出用函数关系做延拓的一般形，
然后按本书后文实际用到的途径归纳方法
（完整构造见第~\ref{ch:gamma_function}~、\ref{ch:riemann_zeta_continuation}~章）。

\paragraph{用语}
正文统一用\textbf{全纯}；与文献中的「解析」同义。区域均指连通开集。


\section{定义与唯一性定理}\label{sec:definition_and_uniqueness}

\begin{definition}{直接解析延拓}{}
设 $D_1$ 和 $D_2$ 是两个区域，且 $D_1\cap D_2\neq\varnothing$。
若 $f_1$ 在 $D_1$ 上全纯，$f_2$ 在 $D_2$ 上全纯，且在 $D_1\cap D_2$ 上有 $f_1=f_2$，
则称 $f_2$ 是 $f_1$ 从 $D_1$ 到 $D_2$ 的\textbf{直接解析延拓}。
\end{definition}

\begin{definition}{解析延拓}{}
设 $f$ 在区域 $D_0$ 上全纯。
如果存在区域 $D\supset D_0$ 和在 $D$ 上全纯的函数 $F$，使得
\[
F|_{D_0} \;=\; f,
\]
则称 $F$ 是 $f$ 的\textbf{解析延拓}（在 $D$ 上的单值全纯延拓）。
\end{definition}

\paragraph{链条式扩张}
延拓可以分步进行：若 $f_0$ 在 $D_0$ 上全纯，并有一串区域 $D_0,D_1,\ldots,D_n$
使相邻两区域相交非空，且每一步都是直接解析延拓，则最终在 $D_n$ 上得到 $f_0$ 的延拓。
（形式化的「圆盘链 / 沿道路延拓」见标准复分析教材；本章只使用上述区域语言。）

\paragraph{单值性说明}
本章只讨论在给定连通区域上的\textbf{单值}全纯延拓。
沿不同路径做直接延拓时，一般可能出现多值（对数、根式等的单值性障碍）。
唯一性定理保证的是：一旦在连通区域 $D$ 上得到一个单值全纯延拓，则它唯一。
$\zeta$ 的标准延拓是 $\mathbb{C}$ 上（仅 $s=1$ 为极点）的单值亚纯函数，无此障碍。

\begin{lemma}{零点孤立性}{}
设 $f$ 在区域 $D$ 内全纯。
若 $z_0\in D$ 满足 $f(z_0)=0$，且 $f$ 在 $D$ 内不恒为零，
则存在 $\delta>0$，使在去心邻域 $0<|z-z_0|<\delta$ 内有 $f(z)\neq 0$。
即：非恒为零的全纯函数的零点是孤立的。
\end{lemma}

\begin{proof}
$f$ 在 $z_0$ 的某邻域内有 Taylor 展开 $f(z)=\sum_{n=0}^\infty a_n(z-z_0)^n$。
因 $f\not\equiv 0$，系数不全为零；又 $f(z_0)=0$，故 $a_0=0$。
设 $a_m$ 为第一个非零系数（$m\ge 1$），则
\[
f(z)=(z-z_0)^m g(z),\qquad
g(z)=\sum_{k=0}^\infty a_{k+m}(z-z_0)^k,
\]
其中 $g$ 在 $z_0$ 附近全纯且 $g(z_0)=a_m\neq 0$。
由连续性，存在 $\delta>0$ 使 $|z-z_0|<\delta$ 时 $g(z)\neq 0$。
从而 $0<|z-z_0|<\delta$ 时 $f(z)\neq 0$。
\end{proof}

\noindent
下文唯一性定理将反复使用本引理。

\begin{theorem}{解析延拓的唯一性定理（恒等定理）}{}
设区域 $D$ 连通，$F$ 与 $G$ 在 $D$ 上全纯。
若存在子集 $E\subset D$，使 $E$ 在 $D$ 内有聚点，且在 $E$ 上 $F=G$，
则在整个 $D$ 上 $F\equiv G$。
\end{theorem}

\begin{proof}
令 $h=F-G$。则 $h$ 在 $D$ 上全纯，且在 $E$ 上 $h=0$。只需证 $h\equiv 0$ 于 $D$。

设 $z_*\in D$ 是 $E$ 的一个聚点。存在点列 $\{z_n\}\subset E\setminus\{z_*\}$ 使 $z_n\to z_*$。
由连续性，$h(z_*)=\lim h(z_n)=0$，故 $z_*$ 也是零点。
于是 $z_*$ 是零点序列的极限，不是孤立零点。
由零点孤立性引理的逆否命题：$h$ 在 $z_*$ 的\textbf{某个}邻域内恒为零
（否则在该邻域内非恒为零的全纯函数不能有非孤立零点）。

为把局部恒零推广到整个连通的 $D$，定义
\begin{align*}
A &= \{ z\in D \mid h^{(k)}(z)=0\ \text{对一切}\ k\ge 0 \},\\
B &= \{ z\in D \mid \text{存在}\ k\ge 0\ \text{使}\ h^{(k)}(z)\neq 0 \}.
\end{align*}
则 $A\cap B=\varnothing$，$A\cup B=D$。

\paragraph{$A$ 是开集}
设 $z_0\in A$。在收敛圆盘 $U(z_0,r)\subset D$ 上，
\[
h(z)=\sum_{k=0}^\infty\frac{h^{(k)}(z_0)}{k!}(z-z_0)^k=0,
\]
故 $h$ 在 $U(z_0,r)$ 上恒为零，从而该圆盘内每点的各阶导数亦为零，即 $U(z_0,r)\subset A$。

\paragraph{$B$ 是开集}
设 $z_1\in B$，则某阶 $h^{(N)}(z_1)\neq 0$。
$h^{(N)}$ 连续，故存在邻域 $U(z_1,\rho)\subset D$ 使其中 $h^{(N)}\neq 0$，从而 $U(z_1,\rho)\subset B$。

$D$ 连通，不能写成两个非空不交开集之并。
已证 $h$ 在 $z_*$ 的某邻域内恒为零，故该邻域上各阶导数为零，即 $z_*\in A$，$A\neq\varnothing$。
因此必有 $B=\varnothing$，即 $A=D$。于是 $h\equiv 0$，即 $F\equiv G$ 于 $D$。
\end{proof}

\begin{corollary}{延拓的唯一性}{}
若用某种方法（级数、围道变形、函数方程等）在更大区域 $D$ 上得到全纯（或亚纯）函数 $F$，
且 $F$ 在原定义域 $D_0$ 上与原函数一致，而 $D$ 连通，
则 $F$ 是原函数在 $D$ 上的\textbf{唯一}单值延拓。
\end{corollary}


\section{用函数关系实现解析延拓的一般原理}\label{sec:fe_continuation_schema}

若已在较小区域 $D_0$ 上知道 $f$，并掌握把「未知点」与「已知点」联系起来的\textbf{函数关系}，
则可用该关系把定义域逐步扩大；再由唯一性，在连通区域上所得单值延拓唯一。

一种常见的范式是：若
\begin{equation}
\label{eq:fe_continuation_schema}
  f(z) \;=\; g(z)\, f\bigl(\Phi(z)\bigr),
\end{equation}
其中 $g$ 在更大区域上已知全纯或亚纯，$\Phi$ 把「未知点」映回 $D_0$（或已延拓区域），
则可用右端定义左端的 $f(z)$。
在 $D_0$ 上与原定义一致时，所得即为原函数在连通区域上的延拓。

\paragraph{极点}
若 $g$ 有极点，则延拓出的 $f$ 一般在对应点有极点。
典型例子：$\Gamma(s)=\Gamma(s+1)/s$ 中 $g(s)=1/s$，在 $s=0,-1,-2,\ldots$ 产生 $\Gamma$ 的极点
（第~\ref{ch:gamma_function}~章）。

常见两类 $\Phi$：
\begin{itemize}
  \item \textbf{平移}：$\Phi(s)=s+1$（如 $\Gamma$ 的递推，见第~\ref{subsec:functional_equation_method}~小节）；
  \item \textbf{反射}：$\Phi(z)=1-z$（$\zeta$ 的函数方程属此类，见第~\ref{ch:riemann_zeta_continuation}~章）。
\end{itemize}

式~\eqref{eq:fe_continuation_schema} 概括\textbf{由函数关系定义延拓}的写法；
下一节按本书实际途径分开叙述，完整构造仍在第~\ref{ch:gamma_function}~、\ref{ch:riemann_zeta_continuation}~章。


\section{本书知识谱系中使用的解析延拓方法}\label{sec:common_methods}

解析延拓的途径很多，例如：Schwarz 反射、Perron / Bromwich 积分、交错级数、
模变换与 Poisson 求和、围道变形、函数方程、Mellin 变换等。
本书\textbf{不必穷尽}；仅按后文需要讨论下列三种
（完整构造见第~\ref{ch:gamma_function}~或~\ref{ch:riemann_zeta_continuation}~章）：

\begin{center}
\renewcommand{\arraystretch}{1.25}
\begin{tabular}{|c|c|c|c|}
\hline
\textbf{途径} & \textbf{本书对象} & \textbf{扩大到的区域（概要）} & \textbf{后文} \\
\hline
交错级数 $\eta$ & $\zeta$ & $\operatorname{Re}s>0$（$s=1$ 极点） & 第~\ref{ch:riemann_zeta_continuation}~章 \\
围道 / Hankel & $\zeta$ & $\mathbb{C}\setminus\{1\}$ 亚纯 & 同上，第一法 \\
函数方程 & $\Gamma$；$\zeta$ & $\Gamma$：全平面（除极点）；$\zeta$：反射对称 & 第~\ref{ch:gamma_function}~、\ref{ch:riemann_zeta_continuation}~章 \\
\hline
\end{tabular}
\end{center}

\subsection{交错级数（$\eta$）}\label{subsec:eta_continuation_method}

思路：另取一个在\textbf{更大半平面}上收敛的级数，建立它与原函数在公共域上的\textbf{代数恒等式}，再用该恒等式定义延拓。

本书实例是 Dirichlet 交错级数
\[
  \eta(s)=\sum_{n=1}^{\infty}\frac{(-1)^{n-1}}{n^s},
\]
它在 $\operatorname{Re}s>0$ 上收敛并给出全纯函数；在 $\operatorname{Re}s>1$ 上有
\[
  \zeta(s)=\frac{\eta(s)}{1-2^{1-s}}.
\]
右端在 $\operatorname{Re}s>0$ 上可用作 $\zeta$ 的定义，从而将 $\zeta$ 延拓到 $\operatorname{Re}s>0$：
\begin{itemize}
  \item $s=1$：分母 $\to 0$ 而 $\eta(1)=\log 2\neq 0$，故为\textbf{真极点}；
  \item 使 $1-2^{1-s}=0$ 的其余点（$2^{1-s}=1$，$s\neq 1$）处 $\eta$ 亦为零，与分母零点相消，为\textbf{可去奇点}。
\end{itemize}
完整推导见第~\ref{ch:riemann_zeta_continuation}~章第~\ref{sec:elementary_extension}~节。
该方法属于\textbf{更好收敛的表示 + 恒等式}的初等延拓。

\subsection{围道与积分表示的变形}\label{subsec:contour_integral_rep}

若函数可写成围道积分（或可变形的路径积分），则往往能把参数移到路径允许的更大区域。
典型形式如
\[
  f(z)=\int_{\gamma}\frac{g(w)}{w-z}\,\mathrm{d}w
\]
（$z$ 不在 $\gamma$ 上时对 $z$ 全纯）；更一般地，将实轴或 Hankel 型路径作连续变形，并控制边贡献。

\textbf{本书主例}：$\zeta$ 的 Riemann \textbf{第一法}——由
$\zeta(s)=\frac1{\Gamma(s)}\int_0^\infty\frac{x^{s-1}}{e^x-1}\,\mathrm{d}x$（$\operatorname{Re}s>1$）
出发，改为 Hankel 围道积分，得到 $s\neq 1$ 上的亚纯延拓并导出与 $\zeta(1-s)$ 的关系。
完整计算见第~\ref{ch:riemann_zeta_continuation}~章第~\ref{sec:riemann_first_method}~节。
（积分表示在收敛半平面上的建立见第~\ref{ch:riemann_zeta_intro}~、\ref{ch:gamma_function}~章。）

\subsection{函数方程}\label{subsec:functional_equation_method}

此处「函数方程」途径即把第~\ref{sec:fe_continuation_schema}~节的
\eqref{eq:fe_continuation_schema} 实现为递推或反射型恒等式，从而扩大定义域。

\paragraph{实例：$\Gamma$ 的递推}
参见第~\ref{ch:gamma_function}~章及图~\ref{fig:gamma-mod-surface}。
取 $D_0=\{\operatorname{Re}s>0\}$，$\Phi(s)=s+1$，$g(s)=1/s$，则
\[
  \Gamma(s) = \frac{\Gamma(s+1)}{s}\qquad (s\neq 0)
\]
为平移型。当 $-1<\operatorname{Re}s\le 0$、$s\neq 0$ 时右端已知，故定义左端；
逐层向左直至全平面（除极点 $0,-1,-2,\ldots$）。构造细节见第~\ref{ch:gamma_function}~章。


\section{解析延拓的障碍：自然边界}\label{sec:natural_boundary}

并非所有函数都能延拓到定义域外。
若收敛圆（或收敛半平面）边界上每一点都是奇点，则无法越过该边界，称其为\textbf{自然边界}。

\begin{example}{自然边界}{}
函数
\[
f(z)=\sum_{n=0}^{\infty} z^{n!}
\]
在 $|z|<1$ 内收敛。单位圆周 $|z|=1$ 是其自然边界：奇点在圆周上稠密，无法解析延拓到圆外任何一点。
（证明提要：在每个单位根方向上，沿半径趋于边界时级数发散，故该边界点为奇点；单位根在圆周上稠密。完整论证见标准教材，此处从略。）
\end{example}

$\zeta(s)$ \textbf{没有}自然边界：它可亚纯延拓到整个复平面，仅 $s=1$ 为极点，是解析延拓的经典成功例子。


\section{解析延拓与黎曼 \texorpdfstring{$\zeta$}{Zeta} 函数}\label{sec:continuation_and_zeta}

黎曼 1859 年对 $\zeta$ 延拓的基本思路可与本章三种途径对照：

\begin{enumerate}
  \item \textbf{级数}：$\zeta(s)=\sum n^{-s}$，$\operatorname{Re}s>1$。
  \item \textbf{积分表示}（与 $\Gamma$ 联系）：
    $\zeta(s)=\frac{1}{\Gamma(s)}\int_0^{\infty}\frac{x^{s-1}}{e^x-1}\,\mathrm{d}x$，$\operatorname{Re}s>1$。
  \item \textbf{围道变形}：改为 Hankel 围道，得到除 $s=1$ 外收敛的积分式，实现亚纯延拓（第一法）。
  \item \textbf{函数方程}：给出 $s$ 与 $1-s$ 两侧的显式关系，揭示临界带对称；
    当延拓在 $\operatorname{Re}s>1$ 上与原级数一致时，由\textbf{唯一性定理}知该单值亚纯延拓唯一。
    （唯一性不来自函数方程本身，而来自恒等定理。）
\end{enumerate}

完整推导见下一章。

\begin{center}
\fbox{\parbox{0.92\textwidth}{
\centering
\vspace{0.2em}
\textbf{第八章：重要定理与核心公式汇总}\\[6pt]
\renewcommand{\arraystretch}{1.6}
\begin{tabular}{lp{9.5cm}}
\hline
\textbf{名称} & \textbf{数学内容 / 公式表达} \\
\hline
解析延拓的唯一性 &
  连通区域上，两全纯函数在有聚点的集合上相等则恒等 \\
单值延拓 &
  本章只讨论连通区域上的单值全纯延拓；$\zeta$ 为全平面单值亚纯 \\
函数关系延拓 &
  $f(z)=g(z)\,f(\Phi(z))$；$g$ 的极点可成为 $f$ 的极点；
  $\Gamma$ 取 $\Phi(s)=s+1$ \\
本书用到的途径 &
  $\eta$ 交错级数；围道 / Hankel；
  函数方程（$\Gamma$ 递推；$\zeta$ 见第~\ref{ch:riemann_zeta_continuation}~章） \\
自然边界 &
  奇点稠密于边界时不可再延拓；$\zeta$ 无自然边界（仅 $s=1$ 极点） \\
\hline
\end{tabular}
\vspace{0.2em}
}}
\end{center}


\section*{本章小结与后续展望}\label{sec:summary_significance}
\addcontentsline{toc}{section}{本章小结与后续展望}

对 $\zeta$ 而言，非平凡零点落在 Dirichlet 级数的收敛域之外：
由 Euler 乘积与函数方程，它们先被限制在闭带 $0\le\operatorname{Re}s\le 1$；
再由 $\operatorname{Re}s=1$ 上无零点得到开临界带 $0<\operatorname{Re}s<1$
（详见第~\ref{ch:zero_distribution}~章）。
因此必须先做解析延拓，才能讨论零点与对称轴 $\operatorname{Re}s=1/2$。

下一章给出 $\zeta$ 的 $\eta$ 初等延拓、Hankel 第一法与 $\vartheta$ 第二法，并推导函数方程；
可与本章三种途径对照阅读。
